\documentclass{amsart}
\usepackage{amsmath,amssymb,amsthm,graphicx}
\newtheorem{theorem}{Theorem}[section]
\newtheorem{lemma}[theorem]{Lemma}
\theoremstyle{definition}
\newtheorem{definition}[theorem]{Definition}
\newtheorem{example}[theorem]{Example}
\title{Balanced colourings of finite sets}
\author{A. Writer}
\begin{abstract}
We prove that every finite set of real numbers can be split into two halves whose sums differ by at most the largest
element, and we show that this bound cannot be improved in general. The proof is a greedy argument, and the paper ends
with two open questions about the same problem for vectors in the plane.
\end{abstract}
\begin{document}
\maketitle

\section{Introduction}\label{sec:intro}

How evenly can a set of numbers be shared between two people? The question goes back to the partition problem studied by
Karp~\cite{karp}, who showed that deciding whether an exact split exists is hard in general. Graham~\cite{graham} gave a
simple rule that is never far from the best split, and Spencer~\cite{spencer} studied the same question for many sets at
once under the name of discrepancy. We prove that the greedy rule always gives a split whose two sums differ by at most the
largest number in the set.

\begin{theorem}\label{thm:main}
For every finite set $A$ of positive real numbers there is a split of $A$ into two parts whose sums differ by at most
$\max A$.
\end{theorem}

In words, no number has to be cut in half: the error of the best split is never more than one element. For example, the
set $\{1,1,1\}$ cannot be split into equal sums, and the best split differs by exactly $1$, which is the largest element,
so the bound of Theorem~\ref{thm:main} is attained. The theorem is proved in Section~\ref{sec:proof}. Compared with the
bound of Graham~\cite{graham}, which is stated for the ratio of the two sums, our bound is new in that it concerns their
difference and holds for every set, not only for large ones. Section~\ref{sec:prelim} fixes the words we use, and every
symbol of the paper is listed in Appendix~\ref{app:notation}, the table at the end.

\section{Preliminaries}\label{sec:prelim}

This section fixes the language of splits and gives a small example of each notion. Throughout, $\mathbb{N}=\{1,2,3,\dots\}$
denotes the positive integers, so the natural numbers start at one.

\begin{definition}\label{def:split}
A \emph{split} of a finite set $A$ is a pair $(B,C)$ of disjoint sets whose union is $A$.
\end{definition}

\begin{example}\label{ex:split}
The set $\{1,2,3\}$ has the split $(\{1,2\},\{3\})$, whose two sums are both equal to three.
\end{example}

\begin{definition}\label{def:gap}
The \emph{gap} of a split $(B,C)$ is the absolute difference of the sum of $B$ and the sum of $C$.
\end{definition}

\begin{example}\label{ex:gap}
The split $(\{1\},\{1,1\})$ of the set $\{1,1,1\}$ has gap one, since its two sums are one and two.
\end{example}

The idea behind everything below is simple: put each number on the lighter side, and the sides can never drift far apart.
Figure~\ref{fig:greedy} shows the rule at work on a small set.

\begin{figure}
\centering
\fbox{greedy rule on the set of three ones}
\caption{The greedy rule places each element on the side whose sum is smaller at that moment.}\label{fig:greedy}
\end{figure}

\section{Proof of the splitting bound}\label{sec:proof}

This section proves Theorem~\ref{thm:main} by following the greedy rule one element at a time.

\begin{lemma}\label{lem:step}
Suppose a split of a set has gap at most $m$ and a number $a$ with $0<a\le m$ is added to the side with the smaller sum.
Then the new split also has gap at most $m$.
\end{lemma}

\begin{proof}
Write $s\le t$ for the two sums before the step, so that $t-s\le m$. After the step the sums are $s+a$ and $t$. If $s+a\le t$
the new gap is $t-s-a$, which is smaller than the old gap. Otherwise the new gap is $s+a-t$, and this is at most $a$, hence at
most $m$. In both cases the gap stays at most $m$, which is what the lemma claims.
\end{proof}

\begin{proof}[Proof of Theorem~\ref{thm:main}]
The plan is to add the elements of $A$ one at a time, each to the lighter side, and to check that the gap never passes
$\max A$. Start with two empty sides, whose gap is zero. Each step adds a number $a\le \max A$, so Lemma~\ref{lem:step} applies
with $m=\max A$ and the gap stays at most $\max A$. When every element has been placed, the two sides form a split of $A$
whose gap is at most $\max A$. This proves Theorem~\ref{thm:main}.
\end{proof}

The hypothesis that the numbers are positive is needed for the rule as stated: with the numbers $5$ and $-5$ the lighter
side changes meaning, and the example in the introduction shows that the bound itself cannot be lowered.

\section{Conclusion}\label{sec:conclusion}

We have shown that a split with gap at most the largest element always exists and that the greedy rule finds one. Two
questions remain open. First, does the same bound hold for vectors in the plane, with the largest length in place of the
largest element? Second, how often does the greedy rule find the best split when the numbers are chosen at random?

\appendix
\section{Table of symbols}\label{app:notation}

\begin{tabular}{lll}
Symbol & Meaning & Where defined \\
$A$ & a finite set of positive real numbers & Theorem~\ref{thm:main} \\
$(B,C)$ & a split of $A$ & Definition~\ref{def:split} \\
$\max A$ & the largest element of $A$ & Theorem~\ref{thm:main} \\
$m$ & a bound on the gap & Lemma~\ref{lem:step} \\
$\mathbb{N}$ & the positive integers & Section~\ref{sec:prelim} \\
\end{tabular}

\begin{thebibliography}{9}
\bibitem{graham} R. L. Graham, Bounds on multiprocessing timing anomalies, SIAM J. Appl. Math. 17 (1969), 416--429.
\bibitem{karp} R. M. Karp, Reducibility among combinatorial problems, in Complexity of Computer Computations, Plenum, New York, 1972, 85--103.
\bibitem{spencer} J. Spencer, Six standard deviations suffice, Trans. Amer. Math. Soc. 289 (1985), 679--706.
\end{thebibliography}
\end{document}
